\input{../tex-inputs/latex-leseansicht-vorspann}

\section[Arthur und Olga Schnitzler an Gustav Schwarzkopf, 27. 8. 1913]{L04514 Arthur und Olga Schnitzler an Gustav Schwarzkopf, 27. 8. 1913}
\nopagebreak\mylabel{L04514v}
\rehead{ }\normalsize\beginnumbering\briefempfaengerindex{Schwarzkopf, Gustav@\textsc{Schwarzkopf, Gustav}!zzzSchnitzler, Olga@\emph{von Olga Schnitzler}!1913-08-272@{27. 8. 1913}|(be}\briefempfaengerindex{Schwarzkopf, Gustav@\textsc{Schwarzkopf, Gustav}!zzzSchnitzler, Arthur@\emph{von Arthur Schnitzler}!1913-08-272@{27. 8. 1913}|(be}
\toendnotes[C]{\smallbreak\pagebreak[2]}
\correspDesc{Versand  durch Arthur Schnitzler, Olga Schnitzler am 27. 8. 1913 in Maloja
\newline{}Übermittlung  am 27. 8. 1913 in Maloja
\newline{}Erhalt  durch Gustav Schwarzkopf im Zeitraum [28. 8. 1913 – 1. 9. 1913?] in Wien}\toendnotes[C]{\smallbreak}
\Standort{DLA, A:Schnitzler, HS.1985.1.1897.}
\physDesc{Bildpostkarte, 192 Zeichen
\newline{}Handschrift Arthur Schnitzler: Bleistift, deutsche Kurrent
\newline{}Handschrift Olga Schnitzler: Bleistift, lateinische Kurrent
\newline{}Versand: 1) Stempel: »\nobreak{}\oindex{Maloja Palace@\textbf{Maloja Palace}|pwk}Majola Palace Majola
                                          (Engadin{[}){]}, 1811 m., Direction: H. Schlangenhauff\pwindex{Schlagenhauff, Hugo *~23.\,8.\,1876 Schwäbisch Hall@\textsc{Schlagenhauff, Hugo} (*~23.\,8.\,1876 Schwäbisch Hall)|pw}\nobreak{}«.   2) Stempel: »\nobreak{}\oindex{Maloja@\textbf{Maloja}|pwk}Maloja (Graubünden), 27. VIII. 13\nobreak{}«. }
\pstart
           \noindent{}{\pb}\textcolor{gray}{\textbf{3005 Ober-Engadin\oindex{Oberengadin@\textbf{Oberengadin}|pw} – Maloja\oindex{Maloja@\textbf{Maloja}|pw} (1811 m) – Kursaal\oindex{Maloja Palace@\textbf{Maloja Palace}|pw} mit Silsersee\oindex{Lej da Segl@\textbf{Lej da Segl}|pw} (1800 m)}}\pend
           \pstart{}{\pb}Hrn \textsc{Gustav Schwarzkopf}\pend{}\pstart{}Wien I\oindex{I., Innere Stadt@\textbf{I., Innere Stadt}|pw}\pend{}\pstart{}\textsc{Tiefer Graben 17}\oindex{Wien@\textbf{Wien}!I., Innere Stadt@\textbf{I., Innere Stadt}!Tiefer Graben 17@\textbf{Tiefer Graben 17}|pw}\pend{}{\bigskip}\vspace{1em}
\pstart
           \centering{}Maloja\oindex{Maloja [Bezirk]@\textbf{Maloja [Bezirk]}|pw}{ }27/8 913\pend
           \vspace{0.5em}
\pstart
           Bitten um ein Wort nach \textsc{St. Moriz\oindex{St. Moritz@\textbf{St. Moritz}|pw}} (Grd \textsc{Hotel}\oindex{Grand Hotel St. Moritz@\textbf{Grand Hotel St. Moritz}|pw}) wo wir wohnen und uns bei ſchönſtem Wetter vortrefflich {\pb}befinden! Ihr
               \spacefill\mbox{A.}\pend
           \selectlanguage{ngerman}\vspace{1em}
\pstart
           \noindent{}{[}hs. O. Schnitzler:{]} Herzlichst Ihre \spacefill\mbox{OlgaS.}\pend
           \selectlanguage{ngerman}\endnumbering\briefempfaengerindex{Schwarzkopf, Gustav@\textsc{Schwarzkopf, Gustav}!zzzSchnitzler, Olga@\emph{von Olga Schnitzler}!1913-08-272@{27. 8. 1913}|)be}\briefempfaengerindex{Schwarzkopf, Gustav@\textsc{Schwarzkopf, Gustav}!zzzSchnitzler, Arthur@\emph{von Arthur Schnitzler}!1913-08-272@{27. 8. 1913}|)be}\mylabel{L04514h}
\begin{anhang}
\end{anhang}\newcommand{\dateiname}{L04514}\newcommand{\titel}{Arthur und Olga Schnitzler an Gustav Schwarzkopf, 27. 8. 1913}\newcommand{\editorInnen}{Herausgegeben von Jahnke, SelmaMüller, Martin Anton}\input{../tex-inputs/latex-leseansicht-abspann}
